% GENERATED FILE. Edit canonical lessons, then run npm run generate:books.
% canonical-source-hash: fnv1a64:f15b3b77facdc102
% canonical-lessons: PA-C65-tamatar, PA-C65-gajar, PA-C65-khand, PA-C65-lun, PA-C65-makkhan

\chapter{Tomato, Carrot, Sugar, Salt, Butter}
\label{ch:pa-tomato-carrot-sugar-salt-butter}
\hlchaptermodality{65}

\begin{chapteropening}
\textbf{By the end of this chapter:} \emph{I can say a tomato, a carrot, sugar, salt and butter.}

Complete the last lesson of chapter 65: i can say a tomato, a carrot, sugar, salt and butter.
\end{chapteropening}

\section[ṭamāṭar]{\pa{ਟਮਾਟਰ} (ṭamāṭar) — a tomato}

\label{lesson:PA-C65-tamatar}



\begin{quote}
Before the new one: say the Punjabi for a banana, then the Punjabi for a potato.
\end{quote}

\subsection*{You'll want to know: \pa{ਟਮਾਟਰ}}
\textbf{\pa{ਟਮਾਟਰ}} — \emph{ṭamāṭar} — \textquotedblleft{}a tomato\textquotedblright{}.

\begin{grammarlens}[title={What you've built}]
One more word for this chapter.
\end{grammarlens}

\subsection*{Your turn}
\begin{itemize}
\raggedright
  \item \emph{Say it:} \emph{ṭamāṭar}
  \item \emph{Say it:} \emph{ṭamāṭar}, once more
  \item \emph{Say it:} say \emph{ṭamāṭar}
  \item \emph{Recall:} say the Punjabi for a banana, then the Punjabi for a potato, then say \emph{ṭamāṭar} again
\end{itemize}

\begin{tcolorbox}[breakable,colback=teal!4,colframe=teal!35!black,title={Before you move on}]
What does \pa{ਟਮਾਟਰ} mean? (\textquotedblleft{}A tomato\textquotedblright{}.) Say it once more.
\end{tcolorbox}
\section[gājar]{\pa{ਗਾਜਰ} (gājar) — a carrot}

\label{lesson:PA-C65-gajar}



\begin{quote}
Before the new one: say the Punjabi for a potato, then the Punjabi for a tomato.
\end{quote}

\subsection*{You'll want to know: \pa{ਗਾਜਰ}}
\textbf{\pa{ਗਾਜਰ}} — \emph{gājar} — \textquotedblleft{}a carrot\textquotedblright{}.

\begin{grammarlens}[title={What you've built}]
One more word for this chapter.
\end{grammarlens}

\subsection*{Your turn}
\begin{itemize}
\raggedright
  \item \emph{Say it:} \emph{gājar}
  \item \emph{Say it:} \emph{gājar}, once more
  \item \emph{Say it:} say \emph{gājar}
  \item \emph{Recall:} say the Punjabi for a potato, then the Punjabi for a tomato, then say \emph{gājar} again
\end{itemize}

\begin{tcolorbox}[breakable,colback=teal!4,colframe=teal!35!black,title={Before you move on}]
What does \pa{ਗਾਜਰ} mean? (\textquotedblleft{}A carrot\textquotedblright{}.) Say it once more.
\end{tcolorbox}
\section[khaṇḍ]{\pa{ਖੰਡ} (khaṇḍ) — sugar}

\label{lesson:PA-C65-khand}



\begin{quote}
Before the new one: say the Punjabi for a tomato, then the Punjabi for a carrot.
\end{quote}

\subsection*{You'll want to know: \pa{ਖੰਡ}}
\textbf{\pa{ਖੰਡ}} — \emph{khaṇḍ} — \textquotedblleft{}sugar\textquotedblright{}.

\begin{grammarlens}[title={What you've built}]
One more word for this chapter.
\end{grammarlens}

\subsection*{Your turn}
\begin{itemize}
\raggedright
  \item \emph{Say it:} \emph{khaṇḍ}
  \item \emph{Say it:} \emph{khaṇḍ}, once more
  \item \emph{Say it:} say \emph{khaṇḍ}
  \item \emph{Recall:} say the Punjabi for a tomato, then the Punjabi for a carrot, then say \emph{khaṇḍ} again
\end{itemize}

\begin{tcolorbox}[breakable,colback=teal!4,colframe=teal!35!black,title={Before you move on}]
What does \pa{ਖੰਡ} mean? (\textquotedblleft{}Sugar\textquotedblright{}.) Say it once more.
\end{tcolorbox}
\section[lūṇ]{\pa{ਲੂਣ} (lūṇ) — salt}

\label{lesson:PA-C65-lun}



\begin{quote}
Before the new one: say the Punjabi for a carrot, then the Punjabi for sugar.
\end{quote}

\subsection*{You'll want to know: \pa{ਲੂਣ}}
\textbf{\pa{ਲੂਣ}} — \emph{lūṇ} — \textquotedblleft{}salt\textquotedblright{}.

\begin{grammarlens}[title={What you've built}]
One more word for this chapter.
\end{grammarlens}

\subsection*{Your turn}
\begin{itemize}
\raggedright
  \item \emph{Say it:} \emph{lūṇ}
  \item \emph{Say it:} \emph{lūṇ}, once more
  \item \emph{Say it:} say \emph{lūṇ}
  \item \emph{Recall:} say the Punjabi for a carrot, then the Punjabi for sugar, then say \emph{lūṇ} again
\end{itemize}

\begin{tcolorbox}[breakable,colback=teal!4,colframe=teal!35!black,title={Before you move on}]
What does \pa{ਲੂਣ} mean? (\textquotedblleft{}Salt\textquotedblright{}.) Say it once more.
\end{tcolorbox}
\section[makkhaṇ]{\pa{ਮੱਖਣ} (makkhaṇ) — butter}

\label{lesson:PA-C65-makkhan}



\begin{quote}
Before the new one: say the Punjabi for sugar, then the Punjabi for salt.
\end{quote}

\subsection*{You'll want to know: \pa{ਮੱਖਣ}}
\textbf{\pa{ਮੱਖਣ}} — \emph{makkhaṇ} — \textquotedblleft{}butter\textquotedblright{}.

\begin{grammarlens}[title={What you've built}]
One more word for this chapter.
\end{grammarlens}

\subsection*{Your turn}
\begin{itemize}
\raggedright
  \item \emph{Say it:} \emph{makkhaṇ}
  \item \emph{Say it:} \emph{makkhaṇ}, once more
  \item \emph{Say it:} say \emph{makkhaṇ}
  \item \emph{Recall:} say the Punjabi for sugar, then the Punjabi for salt, then say \emph{makkhaṇ} again
\end{itemize}

\begin{tcolorbox}[breakable,colback=teal!4,colframe=teal!35!black,title={Before you move on}]
What does \pa{ਮੱਖਣ} mean? (\textquotedblleft{}Butter\textquotedblright{}.) Say it once more.
\end{tcolorbox}
